\section{Cardinais e o número de Hartogs}\label{sec:cardinaisENumeroDeHartogs}
\index{cardinal}
\index{número de Hartogs}

Conforme vimos no Corolário~\ref{cor:finitoCardinal}, para todo conjunto finito $A$, existe um único número natural $n$ tal que $A$ está em bijeção com $n$.
Assim, os números naturais podem ser vistos como representantes canônicos dos conjuntos finitos a menos de bijeção, no que diz respeito à quantidade.

Nesta seção, daremos representantes canônicos também para os conjuntos infinitos bem ordenáveis.
Intuitivamente, dizer que um conjunto é bem ordenável significa que é possível ``contá-lo até o fim utilizando ordinais''.
Neste texto, tais conjuntos serão os únicos para os quais trabalharemos a noção de representação de cardinalidade.
No próximo capítulo, introduziremos o Axioma da Escolha, que implica que todo conjunto é bem ordenável.

Para motivar a nossa definição de cardinal, lembremos como a fizemos para conjuntos finitos.
Associamos, a cada conjunto finito $A$, o número natural $n$ tal que $A$ está em bijeção com $n$.

Se $A$ é bem ordenável, então ele está em bijeção com algum ordinal.
\begin{lemma}\label{lem:bemOrdenavelBijetaComOrdinal}
    Seja $A$ um conjunto bem ordenável.
    Então existe um número ordinal $\alpha$ tal que $A$ está em bijeção com $\alpha$.
\end{lemma}
\begin{proof}
    De fato, seja $<$ uma boa ordem sobre $A$.
    Existe um ordinal $\alpha$ tal que $(A, <)$ é isomorfo a $(\alpha, <)$, pelo Teorema~\ref{thm:boaOrdemEhIsomorfaAOrdinal}, e, portanto, $A$ está em bijeção com $\alpha$.
\end{proof}

Porém, diferentemente do que acontece no caso de ordinais finitos, um conjunto infinito não está em bijeção com um único ordinal.

\begin{example}\label{exa:omegaMaisUmEhEnumeravel}
    $\omega+1$ é enumerável, pois é a união de dois conjuntos enumeráveis: $\omega$ e $\{\omega\}$.
\end{example}
\begin{example}\label{exa:omegaMaisOmegaEhEnumeravel}
    $\omega+\omega$ é enumerável, pois é a união de dois conjuntos enumeráveis: $\omega$ e $\{\omega+n : n\in \omega\}$.
\end{example}

Pelos Exemplos~\ref{exa:omegaMaisUmEhEnumeravel} e \ref{exa:omegaMaisOmegaEhEnumeravel}, se $A$ é infinito e enumerável, então $A$ está em bijeção com $\omega$, com $\omega+1$ e com $\omega+\omega$, dentre muitos outros.
Uma saída natural é associar a $A$ o menor ordinal $\alpha$ tal que $A$ está em bijeção com $\alpha$.

\begin{definition}\label{def:cardinalidadeDeConjuntoBemOrdenavel}\index{cardinalidade}
    Seja $A$ um conjunto bem ordenável.
    A cardinalidade de $A$, ou número de elementos de $A$, é o menor ordinal $\alpha$ tal que $A$ está em bijeção com $\alpha$.
    O cardinal de $A$ é denotado por $|A|$.
\end{definition}

A seguir, vamos tentar explicitar tais objetos de forma mais concreta.
Primeiro, vamos identificar quem são tais ordinais.

\begin{lemma}\label{lem:ordinaisQueSaoCardinalidades}
    Seja $\alpha$ um ordinal.
    Então $\alpha$ é a cardinalidade de algum conjunto bem ordenável se, e somente se, não existe $\beta<\alpha$ tal que $\alpha$ esteja em bijeção com $\beta$.
\end{lemma}
\begin{proof}
    Suponha que não existe $\beta<\alpha$ tal que $\alpha$ esteja em bijeção com $\beta$.
    Temos que $\alpha$ é bem ordenável, e $\alpha$ é o menor ordinal que está em bijeção com $\alpha$, ou seja, $|\alpha|=\alpha$.

    Por outro lado, se $\alpha$ é a cardinalidade de algum conjunto bem ordenável $A$, então $\alpha$ é o menor ordinal que está em bijeção com $A$.
    Logo, não existe $\beta<\alpha$ tal que $\alpha$ esteja em bijeção com $\beta$, pois tal $\beta$ violaria a minimalidade de $\alpha$.
\end{proof}

Com isso, define-se:

\begin{definition}\label{def:cardinal}\index{cardinal}
    Um cardinal é um ordinal $\kappa$ que satisfaz alguma das seguintes condições equivalentes:
    \begin{enumerate}[label=(\roman*)]
    \item $\kappa$ é a cardinalidade de algum conjunto bem ordenável;
    \item não existe $\beta<\kappa$ tal que $\kappa$ esteja em bijeção com $\beta$.
    \end{enumerate}
\end{definition}

\begin{example}
    Todo número natural é um cardinal, pois todo número natural $n$ é um ordinal, e, se $m<n$, temos que $m\subsetneq n$, e, portanto, $m\not\approx n$.
\end{example}
\begin{example}
    $\omega$ é um cardinal, pois é um ordinal, e, se $n<\omega$, então $n\in \omega$, enquanto o item (d) do Corolário~\ref{cor:finitoCardinal} mostra que $\omega$ é infinito.
\end{example}
\begin{example}
    $\omega+\omega$ não é um cardinal, pois $\omega+\omega\approx \omega$.
\end{example}

Um primeiro lema é o seguinte.

\begin{lemma}\label{lem:cardinalEquivNaoBijetaComOrdinalMenor}
    Seja $\alpha$ um ordinal.
    Então $\alpha$ é um cardinal se, e somente se, para todo $\beta<\alpha$, $\alpha$ não está em bijeção com $\beta$.
\end{lemma}
\begin{proof}
    Dado $\beta<\alpha$, temos que $\beta\subseteq \alpha$, e, portanto, $\beta\preceq \alpha$.
    Assim, pelo Teorema de Cantor-Schröder-Bernstein, $\alpha$ está em bijeção com $\beta$ se, e somente se, $\alpha\approx \beta$.
    Desse modo, existe $\beta<\alpha$ tal que $\alpha$ está em bijeção com $\beta$ se, e somente se, existe $\beta<\alpha$ tal que $\alpha\approx \beta$.
\end{proof}

Propriedades elementares dos cardinais são as seguintes.

\begin{lemma}\label{lem:propriedadesElementaresDosCardinais}
    Sejam $A$ e $B$ conjuntos bem ordenáveis.
    Então:
    \begin{enumerate}[label=(\alph*)]
        \item $|A|=|B|$ se, e somente se, $A\approx B$;
        \item $|A|\leq |B|$ se, e somente se, $A\preceq B$;
        \item $|A|<|B|$ se, e somente se, $A\prec B$.
        \item $||A||=|A|$.
    \end{enumerate}
\end{lemma}
\begin{proof}
    (a) Se $|A|=|B|$, então $A\approx |A|=|B|\approx B$.
    Assim, $A\approx B$.

    Reciprocamente, se $A\approx B$, então para todo ordinal $\alpha$, $A\approx \alpha$ se, e somente se, $B\approx \alpha$.
    Assim, $|A|=|B|$.

    (b) Se $|A|\leq |B|$, então $A\preceq |A|\leq |B|\preceq B$, e, portanto, $A\preceq B$.

    Reciprocamente, agimos pela contrapositiva.
    Se $|B|<|A|$, como $|A|$ é um cardinal, $|A|\not\approx |B|$ pelo Lema~\ref{lem:cardinalEquivNaoBijetaComOrdinalMenor}, e, portanto, $|A|\not\preceq |B|$.

    (c) Decorre de (a) e (b).

    (d) Temos que $|A|$ é um cardinal e, portanto, $|A|$ não está em bijeção com nenhum ordinal menor do que ele, de modo que $||A||=|A|$.
\end{proof}

Por enquanto, não provamos a existência de ordinais não enumeráveis, ou, equivalentemente, de boas ordens não enumeráveis.
Assim, todos os cardinais que conhecemos até agora são os naturais e $\omega$.

Com o Axioma da Potência, tal feito se faz possível.

\begin{proposition}\label{prop:teoremaDeHartogs}\index{teorema!de Hartogs}
    Para todo conjunto $A$, existe um ordinal $\alpha$ tal que $A$ não está em bijeção com $\alpha$.
\end{proposition}
\begin{proof}
    Seja $Z=\{(B, R) \in \mathcal P(A) \times \mathcal{P}(A\times A) : R \text{ é uma boa ordem sobre } B\}$.
    
    Pela Substituição, considere o conjunto $\alpha=\{\tp(B, R) : (B, R)\in Z\}$.

    $\alpha$ é um conjunto de ordinais.
    Além disso, $\alpha$ é transitivo.
    Para ver isso, considere $\beta\in \alpha$ e $\xi\in \beta$.
    Seja $(B, R)\in Z$ tal que $\tp(B, R)=\beta$.
    Seja $\phi: \beta \to B$ um isomorfismo.
    Temos que $\phi|_\xi: \xi \to \{b \in B: b R \phi(\xi)\}$ é um isomorfismo, de modo que $\xi\in \alpha$.
    
    Assim, $\alpha$ é um ordinal.

    Temos que $\alpha$ não injeta em $A$: se injetasse, seja $f:\alpha \to A$ uma função injetora, e $B=f[\alpha]$.
    Seja $R$ a relação dada por $f(\xi) R f(\zeta)$ se, e somente se, $\xi<\zeta$.
    Então $f:(\alpha, <)\to (B, R)$ é um isomorfismo, de modo que $\tp(B, R)=\alpha\in \alpha$, um absurdo.
\end{proof}

\begin{definition}\label{def:numeroDeHartogs}\index{número de Hartogs}
    O número de Hartogs de um conjunto $A$, denotado por $h(A)$, é o menor ordinal $\alpha$ tal que $A$ não injeta em $\alpha$.
\end{definition}

\begin{lemma}\label{lem:hartogsEhCardinal}
    Seja $A$ um conjunto.
    Então $h(A)$ é um cardinal.
\end{lemma}
\begin{proof}
    Suponha que não.
    Então existe $\beta<h(A)$ tal que $h(A)$ está em bijeção com $\beta$.
    Pela definição de $h(A)$, $\beta \preceq A$, mas, então, $h(A)\preceq A$, o que é um absurdo.
\end{proof}

\begin{lemma}\label{lem:hartogsEhMenorCardinalMaior}
    Seja $\alpha$ um ordinal.
    Então $h(\alpha)$ é o menor cardinal $\kappa$ tal que $\alpha<\kappa$.
\end{lemma}
\begin{proof}
    Temos que $h(\alpha)\not\preceq \alpha$, logo, $h(\alpha)\neq \alpha$.
    Como $\alpha$ e $h(\alpha)$ são ordinais comparáveis, segue que $\alpha<h(\alpha)$.
    Temos ainda que $h(\alpha)$ é um cardinal.

    Agora seja $\kappa$ um cardinal tal que $\alpha<\kappa$.
    Temos que $\kappa$ não injeta em $\alpha$, e, portanto, $h(\alpha)\leq \kappa$.
\end{proof}

\begin{definition}\label{def:sucessorCardinal}\index{sucessor cardinal}
    Dado um ordinal $\alpha$, o número de Hartogs de $\alpha$ é denotado por $\alpha^+$.
\end{definition}

Por fim, temos que a união de uma coleção de cardinais é um cardinal.

\begin{lemma}\label{lem:uniaoDeCardinaisEhCardinal}
    Seja $A$ uma coleção de cardinais.
    Então $\bigcup A$ é um cardinal.
\end{lemma}
\begin{proof}
    Sabemos que a união de ordinais é um ordinal.
    Para ver que $\bigcup A$ é um cardinal, tome $\beta<\bigcup A$.
    Seja $\kappa\in A$ tal que $\beta<\kappa$.
    Como $\kappa$ é um cardinal, $\kappa$ não injeta em $\beta$, e, portanto, $\bigcup A$, que contém $\kappa$, não injeta em $\beta$.
\end{proof}

Finalizamos com um lema que restringe a estrutura dos cardinais.

\begin{lemma}\label{lem:cardinalInfinitoEhOrdinalLimite}
Seja $\kappa$ um cardinal infinito.
Então $\kappa$ é um ordinal limite.
\end{lemma}
\begin{proof}
    Por definição, $\kappa$ é um ordinal.
    Como ele é infinito, $\kappa\neq 0$.
    Suponha por absurdo que $\kappa=\alpha+1$ para algum ordinal $\alpha$.
    Como $\kappa$ é infinito e $\kappa=\alpha\cup\{\alpha\}$, temos que $\alpha$ é infinito.
    Assim, $\omega\subseteq \alpha$.
    Seja $f:\kappa \to \alpha$ dada por:
    \[
    f(\xi) = \begin{cases}
        \xi+1, & \text{se } \xi<\omega\\
        \xi, & \text{se } \omega \leq \xi < \alpha\\
        0, & \text{se } \xi=\alpha
    \end{cases}
    \]
    É fácil verificar que $f$ é uma função bijetora, de modo que $\kappa$ não é um cardinal.
\end{proof}

\subsection*{Exercícios}
Nos exercícios dessa seção, NÃO assuma o Axioma da Escolha (este ainda não foi introduzido!).
\begin{exercise}
    Mostre que o ordinal $\alpha$ construído na demonstração da Proposição~\ref{prop:teoremaDeHartogs} é o número de Hartogs de $A$.
\end{exercise}
\begin{exercise}
    Seja $A$ um conjunto.
    Mostre que $h(A)$ é o menor ordinal tal que $A$ não injeta em $h(A)$.
\end{exercise}
\begin{exercise}
    Seja $\kappa$ um cardinal.
    Mostre que não existe um cardinal $\lambda$ tal que $\kappa<\lambda<\kappa^+$.

    É possível que existam ordinais $\alpha$ tais que $\kappa<\alpha<\kappa^+$?
    Em que casos?
\end{exercise}
\begin{exercise}
Mostre que se $\alpha$ é um ordinal infinito, então $\alpha+1\approx \alpha$.
\end{exercise}
\begin{exercise}
    Mostre que se $A\subseteq B$ e $B$ é bem ordenável, então $|A|\leq |B|$.
\end{exercise}
